\documentclass[10pt,a4paper]{amsart}
\usepackage[margin=19mm]{geometry}
\usepackage{amsmath,amssymb,mathtools}
\usepackage[scr=rsfs,cal=euler]{mathalfa}
\usepackage{xfrac}
\usepackage[unicode,hidelinks,bookmarks=false]{hyperref}
\newcommand{\ie}{i.e.\ }
\newcommand{\cf}{cf.\ }
\newcommand{\resp}{respectively\ }
\newcommand{\Subsection}{\subsection}
\providecommand{\nobreakdash}{\nobreak\hbox{-}}
\setlength{\emergencystretch}{2em}
\multlinegap0pt
\usepackage{amssymb}
\usepackage[shortlabels]{enumitem}
\usepackage{mathtools}
\usepackage{tikz}
\newtheorem{lem}{Lemma}
\title{On General Linear Degenerate Elliptic PDE Systems}
\author{Nikos Katzourakis, Frederick Temple}
\date{Source: arXiv:2607.22888v1 (CC BY 4.0)}
\begin{document}
\maketitle

\noindent\emph{Source and licence.} On General Linear Degenerate Elliptic PDE Systems. Version arXiv:2607.22888v1 (CC BY 4.0), distributed by arXiv under the Creative Commons Attribution 4.0 International licence, http://creativecommons.org/licenses/by/4.0/, as recorded in the arXiv licence metadata for this exact version and shown on the abstract page https://arxiv.org/abs/2607.22888v1. Full licence notice: \texttt{licenses/arxiv-2607.22888-CC-BY-4.0.txt}.\par\smallskip

\noindent\textit{Editorial scope.} Counted unit: the article's lem lem:1 (source0.tex lines 499--504). The article's complete original proof of it is delivered with it (source0.tex lines 506--533, one \texttt{proof} environment of 28 lines). No mathematics is written, repaired or re-derived: the statement and the proof are the article's own lines, in the article's own notation, and no retained line is substituted or re-flowed.  No further source-local context is needed: every cross-reference of every retained line resolves inside the two delivered spans.  Nothing was omitted from any retained span, so the package carries no omission record.  No retained line contains a citation, so the wrapper carries no bibliography and no bibliography entry.  No retained line contains a figure environment, an image-inclusion command or drawing code, so no image is delivered and the package declares no asset dependency.  Editorial formatting only, in addition: an AMS \texttt{amsart} preamble that restates the article's own declarations and macros used by the retained lines, namely the environments \texttt{lem} on separate counters, together with the standard packages those lines need.  The retained environments print in this document as Lemma 1 in order of appearance, whereas the article prints the same items with the numbering of its own full document.  The complete native source of the article as distributed in the arXiv e-print for this exact version is bundled unchanged as \texttt{sources/205978/source0.tex}.
\begin{lem} \label{lem:1}
Let $\mathbf{B}: \mathbb{R}^{Nn} \longrightarrow \mathbb{R}^{N}$ be a linear map satisfying the symmetry condition $\mathbf{B}_{\alpha \beta i} = \mathbf{B}_{\beta \alpha i}$. Then, for any $u \in W^{1,2}_0(\Omega, \mathbb{R}^N)$, we have
\begin{equation*}
\int_{\Omega} \sum_{\alpha, \beta, i}\mathbf{B}_{\alpha \beta i} \mathrm{D}_i u_{\beta} u_{\alpha} = 0.
\end{equation*}
\end{lem}

\begin{proof}[Proof of Lemma \ref{lem:1}]
Fix $u \in W^{1,2}_0(\Omega, \mathbb{R}^N)$. Note first that we can rewrite the term $(\mathrm{D}_i u_\beta)u_\alpha$ as 
\[
\mathrm{D}_i(u_\alpha u_\beta) - (\mathrm{D}_i u_\alpha)u_\beta. 
\]
Therefore, we have the idenity
\begin{equation}\label{eq:5new}
\int_{\Omega} \sum_{\alpha, \beta, i}\mathbf{B}_{\alpha \beta i} (\mathrm{D}_i u_{\beta}) u_{\alpha} = \int_{\Omega} \sum_{\alpha, \beta, i}\mathbf{B}_{\alpha \beta i} \mathrm{D}_i(u_{\alpha} u_{\beta}) - \int_{\Omega} \sum_{\alpha, \beta, i}\mathbf{B}_{\alpha \beta i} (\mathrm{D}_i u_{\alpha}) u_{\beta}.
\end{equation}
Integrating by parts the first term on the right-hand side of equation \eqref{eq:5new} leads to

\begin{equation}
\int_{\Omega} \sum_{\alpha, \beta, i}\mathbf{B}_{\alpha \beta i} \mathrm{D}_i(u_{\alpha} u_{\beta}) = \int_{\partial \Omega} \sum_{\alpha, \beta, i}\mathbf{B}_{\alpha \beta i}(u_{\alpha} u_{\beta})\nu_i = 0,
\end{equation}
where the last equality owes to that $u \in W^{1,2}_0(\Omega, \mathbb{R}^N)$. Therefore, we arrive at
\begin{equation} \label{eq:4new}
\int_{\Omega} \sum_{\alpha, \beta, i}\mathbf{B}_{\alpha \beta i} (\mathrm{D}_i u_{\beta}) u_{\alpha} = - \int_{\Omega} \sum_{\alpha, \beta, i}\mathbf{B}_{\alpha \beta i} (\mathrm{D}_i u_{\alpha}) u_{\beta}.
\end{equation}
Under the assumption that $\mathbf{B}_{\alpha \beta i} = \mathbf{B}_{\beta \alpha i}$, we obtain
\[
\int_{\Omega} \sum_{\alpha, \beta, i}\mathbf{B}_{\alpha \beta i} (\mathrm{D}_i u_{\beta}) u_{\alpha} = \int_{\Omega} \sum_{\alpha, \beta, i}\mathbf{B}_{\beta \alpha i} (\mathrm{D}_i u_{\alpha}) u_{\beta}= \int_{\Omega} \sum_{\alpha, \beta, i}\mathbf{B}_{\alpha \beta i} (\mathrm{D}_i u_{\alpha}) u_{\beta},
\]
where the first equality is due to the relabelling of dummy indices. Therefore, by equation \eqref{eq:4new} and the above, we deduce
\begin{equation*}
\int_{\Omega} \sum_{\alpha, \beta, i}\mathbf{B}_{\alpha \beta i} (\mathrm{D}_i u_{\beta}) u_{\alpha} = 0.
\end{equation*}
The conclusion ensues.
\end{proof}

\end{document}
